\documentclass[a4paper]{article}
\usepackage{amsfonts}
\usepackage{amsmath}
\usepackage{amssymb}
\usepackage{graphicx}     
\usepackage{cancel}

% Command definities van frequent gebruikte commandos
\newcommand{\N}{\mathbb{N}}
\newcommand{\Q}{\mathbb{Q}}
\newcommand{\R}{\mathbb{R}}
\newcommand{\C}{\mathbb{C}}
\newcommand{\Bgsin}{\operatorname{Bgsin}}
\newcommand{\Bgcos}{\operatorname{Bgcos}}
\newcommand{\Bgtan}{\operatorname{Bgtan}}
\newcommand{\cis}{\operatorname{cis}}
\newcommand{\Limit}{\lim\limits_}
\newcommand{\Int}{\displaystyle\int}
\newcommand{\Dint}{\displaystyle\int\limits}

\begin{document}
  
\noindent \large Auteur: Wout Vanderborght\\
\noindent \large Studierichting: Informatica\\
\noindent \large Oefeningenreeks 6C nr. 4.10\\

\medskip

\normalsize

$ x_n \dfrac{6n^3+4}{(n+1)(n+2)(1-2n)} $ \\

$ \Limit{n \rightarrow +\infty} x_n = \dfrac{6 \cdot \cancel{n^3}}{2 \cdot \cancel{n^3}} = 3 $ \\

$ \mathop{\operatorname{adh}}\limits_{n \rightarrow +\infty} x_n = \{3\} $ \\

$ \liminf\limits_{n \rightarrow +\infty} x_n = 3 $ \\

$ \limsup\limits_{n \rightarrow +\infty} x_n = 3 $ \\

\begin{thebibliography}{999}
%\bibitem{a} Referentie
\end{thebibliography}
\end{document} 
